%!TeX program = lualatex
\renewcommand{\thelektion}{L04 \textperiodcentered{} Rechnen im Binärsystem \& Bits/Bytes}

\begin{frame}\lessontitle{Lektion 04}{Rechnen im Binärsystem \& Bits/Bytes}{Schriftliche Addition, Überlauf und die kleinste Informationseinheit \textperiodcentered{} Skript, Kapitel 2/3}\end{frame}

\begin{frame}\FDUtitle{Rechnen mit Zahlen}
Addieren im Dezimalsystem ($49035_{10}$ + $718293_{10}$):
\[
\begin{array}{ccccccc}
  & 0 & 4 & 9 & 0 & 3 & 5_{10} \\
+ & \underset{}{7} & \underset{1}{1} & \underset{}{8} & \underset{1}{2} & \underset{}{9} & \underset{}{3_{10}} \\
\hline
 & 7 & 6 & 7 & 3 & 2 & 8_{10} \\\hline\hline
\end{array}
\]
\begin{center}Genau gleich funktioniert das Ganze im Binärsystem!\end{center}
\end{frame}

\begin{frame}\FDUtitle{Rechnen im Binärsystem}
Nur 5 Regeln:
\begin{itemize}
\item 0 + 0 = 0
\item 0 + 1 = 1
\item 1 + 0 = 1
\item 1 + 1 = 0, Übertrag 1
\item 1 + 1 + 1 = 1, Übertrag 1
\end{itemize}

Beispiel ($101011_2$ + $111010_2$):
\[
\begin{array}{ccccccc}
	& 1 & 0 & 1 & 0 & 1 & 1_2 \\
	+ & \underset{1}{1} & \underset{1}{1} & \underset{}{1} & \underset{1}{0} & \underset{}{1} & \underset{}{0_2} \\
	\hline
	1 & 1 & 0 & 0 & 1 & 0 & 1_2 \\\hline\hline
\end{array}
\]
\end{frame}

\begin{frame}\FDUtitle{Was passiert bei Überlauf?}
\begin{center}
Rechnen Sie $11110001_2$ + $00001111_2$ in einer \textbf{8-Bit-Umgebung}. Was fällt auf?
\end{center}
\only<2->{
\begin{center}
Mathematisch korrekt wäre $100000000_2$ -- das sind aber \textbf{9} Bits!
\end{center}
\only<3->{\begin{center}\color{FDUteal}\bfseries In einer starren 8-Bit-Umgebung geht die führende 1 verloren $\rightarrow$ \tib{Overflow} (Überlauf).\end{center}}
}
\end{frame}

\begin{frame}{Auftrag}{Kapitel 2: Rechnen im Binärsystem}
\aufgabebox{\textbf{Aufgabe 2.11, 2.12, 2.14}: Binäre Additionen ($10001010_2 + 11101_2$) \& grösste 4-stellige Binärzahl}\\[3mm]
\challengebox{\textbf{Aufgabe 2.13 (Challenge)}: Darstellung der grössten $n$-stelligen Binärzahl $z$}
\end{frame}

{\setbeamertemplate{footline}{\darkfootline}
\begin{frame}\sectionframe[FDUblue]{Bits \& Bytes}{Wie zählt man mit Nullen und Einsen?}
\end{frame}}

\begin{frame}\FDUtitle{Bits \& Bytes}
\begin{itemize}[<+->]
	\item \textbf{bit} = \enquote{\textit{\underline{bi}nary digi\underline{t}}}, \enquote{binäre Ziffer} (eine 0 oder eine 1)
	\item \textbf{byte} = \enquote{Bissen} (von engl.\ \textit{bite} = Bissen, aus 8 bits) \emoji{yum}
	\item Mit einem Byte lassen sich $2^8=256$ verschiedene Werte darstellen
\end{itemize}
\end{frame}

\begin{frame}\FDUtitle{Führende Nullen}
Die binäre Repräsentation der Dezimalzahl 22:
\[00010110_2 = 22_{10}\]
\begin{center}
Bei Binärzahlen werden -- anders als bei Dezimalzahlen -- üblicherweise \textbf{alle} Bits ausgeschrieben.
\end{center}
\vspace{3mm}
\begin{center}
\textbf{Frage}: Wie stellt man $42_{10}$ mit einem Byte dar?
\end{center}
\only<2->{\begin{center}\Large $42_{10} = 00101010_2$\end{center}}
\end{frame}

\begin{frame}\FDUtitle{Grössenordnungen}
\begin{table}[H]
	\footnotesize
	\centering
	\begin{tabular}{@{}l|rl@{}}
		\toprule
		\textbf{Einheit} & \textbf{Grösse} & Beispiel \\ \midrule
		KB (Kilobyte) & $10^3$ B & eine Text-Datei \\
		MB (Megabyte) & $10^6$ B & eine Musik-Datei \\
		GB (Gigabyte) & $10^9$ B & eine Video-Datei \\
		TB (Terabyte) & $10^{12}$ B & kleiner Firmen-Server \\
		\bottomrule
	\end{tabular}
\end{table}
\end{frame}

\begin{frame}{Auftrag}{Kapitel 3: Bits und Bytes}
\aufgabebox{\textbf{Aufgabe 3.1}: Darstellung von $42_{10}$ mit einem Byte (inkl. führende Nullen)}\\[3mm]
\aufgabebox{\textbf{Aufgabe 3.2}: Grösste Dezimalzahl mit einem Byte (= 8 Bits) bestimmen}
\end{frame}

